\section{Task 5: \textit{NormWear on the Cloud}}

All experiments ran on a Google Cloud Compute Engine VM (e2-standard-4: 4 vCPU, 16 GB RAM, no GPU) using PyTorch on CPU. We used only the three \texttt{total\_acc} channels of UCI HAR because the ADXL345 on our ESP32 has no gyroscope, so the same input format carries over to Task 5.3.

\subsection{Task 5.1: Zero-Shot Inference}

\textbf{Inputs to NormWear:}
\begin{itemize}
    \item \textbf{Task description:} ``What is the current activity?''
    \item \textbf{Answer options:} ``The person is walking.'', ``The person is walking upstairs.'', ``The person is walking downstairs.'', ``The person is sitting.'', ``The person is standing.'', ``The person is laying.''
    \item \textbf{Signal:} raw UCI HAR test windows in g, with shape $[B, 3, 128]$ at 50 Hz and no normalization.
\end{itemize}

The text is encoded by NormWear's text encoder, and the signal by the NormWear encoder and MSiTF fusion module (\texttt{signal\_encode}). The predicted class is the option closest to the signal embedding. Because \texttt{scipy.signal.cwt} has been removed from recent SciPy versions, we enabled the model's PyTorch wavelet option.

\begin{table}[h]
\centering
\begin{tabular}{lr}
\toprule
Metric & Value \\
\midrule
Test windows & 2947 \\
Zero-shot accuracy & 11.27\% \\
Runtime & 3622 s ($\approx$ 60 min) \\
Hardware & CPU \\
\bottomrule
\end{tabular}
\caption{Zero-shot results on the UCI HAR test set}
\end{table}

The model predicted only two classes: WALKING\_DOWNSTAIRS (1257 windows) and SITTING (1690 windows). Its accuracy is below uniform random guessing (16.7\%) and below the handout's reference of about 16\%. The runtime is also longer than the handout's estimate of 10--20 minutes. The VM was hosting our Flask server during the run, so this time is not a clean benchmark.

\subsection{Task 5.2: Few-Shot Linear Probe}

We kept the NormWear encoder frozen and extracted embeddings with \texttt{get\_embedding}, averaging its $[B, 3, P, 768]$ output over patches and channels to obtain one 768-dimensional vector per window. We encoded 10 randomly selected training windows per class (60 total) and all 2947 test windows, which took 3386 s on CPU. We then trained one \texttt{nn.Linear(768, 6)} layer with cross-entropy loss and Adam (learning rate $10^{-3}$) for 100 epochs.

\begin{table}[h]
\centering
\begin{tabular}{lr}
\toprule
Model & Test accuracy \\
\midrule
Before: zero-shot NormWear (Task 5.1) & 11.27\% \\
Untrained linear head (random initialization) & 22.46\% \\
After: trained linear head & 62.61\% \\
\bottomrule
\end{tabular}
\caption{Test accuracy before and after training the linear head}
\end{table}

Training the linear head improved accuracy by 51.3 percentage points over the zero-shot result, using only 60 labeled windows.

\subsection{Task 5.3: End-to-End Pipeline}

\begin{enumerate}
    \item The ESP32 samples the ADXL345 at 50 Hz for 128 samples per axis (about 2.56 s), converting readings to g.
    \item It sends the $3 \times 128$ window as JSON in an HTTP POST to the Flask \texttt{/classify} endpoint on the VM.
    \item The server converts the window to a $[1, 3, 128]$ tensor and passes it through the frozen NormWear encoder to obtain a $[1, 768]$ embedding. The encoder and linear head are loaded once at server startup.
    \item The trained linear head and softmax produce class probabilities, and the most probable class is the prediction.
    \item The server returns the label and confidence to the ESP32, which displays them on the OLED. Each request is also appended to \texttt{predictions.jsonl}.
\end{enumerate}

Table~\ref{tab:esp32-example} shows one logged run. The board was at rest and not worn on the body, so this example demonstrates that the pipeline works rather than recognition accuracy.

\begin{table}[h]
\centering
\begin{tabular}{lr}
\toprule
Field & Value \\
\midrule
Input & $3 \times 128$ ADXL345 window, 50 Hz \\
Recording time & 2542 ms \\
Server inference time & 1248 ms \\
\textbf{Prediction} & \textbf{STANDING (95.6\%)} \\
\midrule
P(STANDING) & 0.956 \\
P(SITTING) & 0.023 \\
P(WALKING\_DOWNSTAIRS) & 0.010 \\
P(WALKING) & 0.006 \\
P(LAYING) & 0.004 \\
P(WALKING\_UPSTAIRS) & 0.001 \\
\bottomrule
\end{tabular}
\caption{Example ESP32 input and prediction returned by the cloud server}
\label{tab:esp32-example}
\end{table}

% Optional: add a photo of the OLED or a screenshot of the /predictions page.
% \begin{figure}[h]
%     \centering
%     \includegraphics[width=0.7\linewidth]{oled_prediction.jpg}
%     \caption{Predicted activity displayed on the ESP32 OLED}
% \end{figure}
